\section{A uniform lattice local limit estimate}
\label{app:lattice-llt}

\paragraph{Lemma.}
Let $X$ have a log-concave probability mass function with interval
support in $\mathbb Z$, mean $\mu$, and finite variance $\sigma^2=\mathrm{Var}(X)\ge1$.
If $X_1,\ldots,X_n$ are independent copies and $n\mu\in\mathbb Z$, then
\[
\mathbb P(X_1+\cdots+X_n=n\mu)
=\frac{1+O(n^{-1/2})}{\sqrt{2\pi n}\,\sigma},
\]
with a numerical implied constant and no upper bound on $\sigma$.
If the law belongs to an exponential tilt family, then
$|d\sigma/d\mu|\le C$ wherever $\sigma\ge1$ and the tilt is interior.

\paragraph{Modal mass and moments.}
Write $p_j=\mathbb P(X=j)$, choose a mode $m$ and put $b=p_m$.
Chebyshev's inequality gives $b\ge c/\sigma$ by placing at least
$3/4$ of the mass in an interval containing at most $4\sigma+2$
integers. On either side of the mode let $J$ be the first distance
where the mass is at most $b/2$, extending by zero off the support.
Since the total mass is at most one, $J\le 2/b+1\le3/b$.
Concavity of the log mass gives $p_{m\pm j}\le b\,2^{-j/J}$ for
$j\ge J$ (immediate if the support ends at $J$), and the modal bound gives $p_{m\pm j}\le2b\,2^{-j/J}$
for $0\le j<J$. Thus $p_{m\pm j}\le2b\,2^{-bj/3}$. Summing $j^2$
against this bound gives $\sigma^2\le C/b^2$ and hence $b\le C/\sigma$.
The same sum with $j^3$ gives $\mathbb E|X-m|^3\le C\sigma^3$,
and centering preserves this bound by Jensen's inequality and the
triangle inequality. In an interior exponential tilt with parameter
$\eta$, the cumulant identities
\[
\frac{d\mu}{d\eta}=\sigma^2,\qquad
\frac{d\sigma^2}{d\eta}=\mathbb E(X-\mu)^3
\]
give $|d\sigma/d\mu|\le C$.

\paragraph{Fourier bounds.}
Let $\phi(t)=\mathbb E e^{it(X-\mu)}$ for $|t|\le\pi$. For $\sigma$
above a fixed numerical threshold, consider the sites in
$[\mu-2\sigma,\mu+2\sigma]$ of mass at least $1/(24\sigma)$.
At most $6\sigma$ sites are in this interval, so the retained sites
have mass at least $1/2$. Since $b\le C/\sigma$, there are at least
$c\sigma$ retained sites, and they are contiguous by log-concavity.
They give a block of length $M\asymp\sigma$, with $M\ge16$ and mass
at least $c/M$ at each site. Restricting
\[
1-|\phi(t)|^2=\sum_{i,j}p_ip_j\bigl(1-\cos((i-j)t)\bigr)
\]
to this block bounds it below by a numerical multiple of
$1-|U_M(t)|^2$, where $U_M(t)=M^{-1}\sum_{j=0}^{M-1}e^{ijt}$.
If $2\pi/M\le|t|\le\pi$, the sine quotient gives
$|U_M(t)|\le\pi/(M|t|)\le1/2$. If $|t|\le2\pi/M$, retain the pairs
whose distances lie between $M/8$ and $M/4$. There are at least
$cM^2$ such pairs, and $1-\cos((i-j)t)\ge cM^2t^2$ on them. This
gives
\[
|\phi(t)|\le \exp\bigl(-c\min(\sigma^2t^2,1)\bigr).
\]

For $1\le\sigma$ below the fixed threshold, the modal lower bound
makes $b$ bounded away from zero and the crossing distances $J$
bounded. On either side of the mode, up to distance $J$ the masses
are at most that of the immediate neighbor on that side. Beyond
$J$, the same neighbor mass bounds the tail after multiplication
by $2^{-(j-J)/J}$. Thus $\mathbb E(X-m)^2\le C(p_{m-1}+p_{m+1})$.
Since this expectation is at least $\sigma^2\ge1$, one neighboring
mass is bounded below. Pairing it with the modal mass in the
cosine sum proves the same Fourier bound for bounded $\sigma$.
Finally, padded unimodality gives total variation $2b$, so summation
by parts yields
\[
|\phi(t)|\le \frac{2b}{|1-e^{it}|}\le\frac{C}{\sigma|t|}
\qquad(0<|t|\le\pi).
\]

\paragraph{Inversion.}
Fix a numerical $0<a<1$. For $|t|\le a/\sigma$, Taylor's formula
and the third-moment bound give
$|\phi(t)-e^{-\sigma^2t^2/2}|\le C\sigma^3|t|^3$.
The Fourier bound above gives Gaussian decay on this interval.
Telescoping the powers of order $n$ therefore bounds their difference
by $C n\sigma^3|t|^3 e^{-c(n-1)\sigma^2t^2}$, whose integral is
$O((n\sigma)^{-1})$ for $n\ge2$.
The Gaussian integral outside this interval is exponentially small
relative to $(\sqrt n\,\sigma)^{-1}$. On
$a/\sigma\le |t|\le A/\sigma$, with $A$ a fixed large constant,
the original Fourier bound gives an integral $O(e^{-cn}/\sigma)$.
On the rest, integrate $(C/(\sigma |t|))^n$ from $A/\sigma$ to
infinity, choosing $A>2C$. This gives at most
$2A(C/A)^n/(\sigma(n-1))$. The intervals are intersected with
$[-\pi,\pi]$ and empty parts are omitted.
Since $n\mu$ is integral, Fourier inversion gives the central mass
as $(2\pi)^{-1}\int_{-\pi}^{\pi}\phi(t)^n\,dt$. Comparison with
the full Gaussian integral proves the lemma.
